% TOP_PROBLEM: {"id": 2302032, "title": "Research Problems in Function Theory — Problem 2.32", "category_id": 8, "category": {"id": 8, "name": "analysis", "display_name": "Analysis", "description": "Limits, continuity, calculus, and function theory.", "slug": "analysis", "order_index": 8, "created_at": "2026-07-31T15:26:25.671Z"}, "status": "open"}
% FIRST_POSTED: null
% ENTRY_KIND: statement_only
% Imported from: batch11.tex; UnsolvedMath ID 2302032
% Compile twice: lualatex 2302032.tex
\documentclass[11pt]{article}
\usepackage{amsmath,amssymb,mathtools}
\usepackage{fontspec}
\usepackage{unicode-math}
\setmainfont{Latin Modern Roman}
\setmathfont{Latin Modern Math}
\usepackage{xurl}
\usepackage[unicode,colorlinks=true,linkcolor=blue,urlcolor=blue,pdftitle={UnsolvedMath Problem 2302032}]{hyperref}
\newcommand{\sref}[2]{\hyperlink{source:#1:#2}{[S#2]}}
\newcommand{\cataloguescope}[1]{\paragraph{Mathematical statement.}#1}
\begin{document}
% UnsolvedMath ID 2302032; source catalogue code AMR-022-2032
\setcounter{section}{2302031}
\section{Research Problems in Function Theory — Problem 2.32}\label{2302032}
{\small\sffamily UnsolvedMath ID 2302032 · Analysis}
\subsection{Definitions and mathematical statement}

\paragraph{Shared notation.}$\mathbb N=\{1,2,\ldots\}$, and $\mathbb Z,\mathbb Q,\mathbb R,\mathbb C$ denote the integers, rationals, reals and complex numbers. Any other conventions are specified locally.


Let $f(z)=\sum_{n=0}^\infty a_nz^n$ be transcendental entire with $a_n\ge0$, and for $x>0$ set $p_n(x)=a_nx^n/f(x)$, so that $\sum_{n\ge0}p_n(x)=1$.
The integral constraints below imply $a_n>0$ for every $n$; integrals are improper integrals on the positive real axis.

\paragraph{Statement recorded in the supplied source.}
Do the identities
\[\int_0^\infty \frac{a_nx^n}{f(x)}\,dx=1\quad(n=0,1,2,\ldots),\]
or equivalently
\[\int_0^\infty \frac{f(tx)}{f(x)}\,dx=\frac1{1-t}\quad(0<t<1),\]
force $f(z)=ce^z$ for some $c>0$? \sref{2302032}{2}
Update 2.32 explicitly gives the scale-corrected uniqueness question $f(z)=ce^z$, rather than the original literal exclusion of only $e^z$, and records Miles and Williamson's affirmative theorem. \sref{2302032}{2}
\subsection{Short English statement}
Do the prescribed integral identities uniquely determine the exponential entire function up to a positive constant factor?
\subsection{Sources}
\begin{description}
\item[\hypertarget{source:2302032:1}{S1}] ULAM.AI and the UnsolvedMath Contributors, \emph{UnsolvedMath: A Curated Collection of Open Mathematics Problems}, record 2302032 (AMR-022-2032). CC BY 4.0. \url{https://huggingface.co/datasets/ulamai/UnsolvedMath}.
\item[\hypertarget{source:2302032:2}{S2}] W. K. Hayman and E. F. Lingham, \emph{Research Problems in Function Theory} (2018), Chapter 2, Problem 2.32 and Update 2.32. \url{https://arxiv.org/abs/1809.07200}.
\end{description}
\label{end:2302032}
\end{document}
